\uuid{nIbu}
\exo7id{1293}
\titre{exo7 1293}
\auteur{gourio}
\organisation{exo7}
\datecreate{2001-09-01}
\isIndication{false}
\isCorrection{false}
\chapitre{Calcul d'intégrales}
\sousChapitre{Intégrale impropre}
\cpge{mpsi-mp2i.nombres-reels-suites.d, pcsi.nombres-reels-suites.c, ptsi.nombres-reels-suites.c, mpsi-mp2i.analyse-asymptotique.a, pcsi.analyse-asymptotique.a, ptsi.analyse-asymptotique.a, mpsi-mp2i.procedes-sommatoires-discrets.a, pcsi.series-numeriques.a, ptsi.series-numeriques.a, bcpst1.suites-reelles, mp-mpi.series-numeriques-vectorielles.b, pc.series-numeriques, psi.series-numeriques.a, pt.series-numeriques.a, bcpst2.series-reelles, bcpst2.integrales-generalisees}
\module{Analyse}
\niveau{L1}
\difficulte{}

\contenu{
\texte{
Soit $f\in C([1,\infty [,{\Rr}^{+})$ d\'{e}croissante, on pose :
$$x_{n}=\sum_{k=1}^{n}f(k)-\int_{1}^{n}f(t)dt. $$
}
\begin{enumerate}
    \item \question{Montrer que la suite $\left( x_{n}\right) _{n\in \Nn} $ converge.}
    \item \question{Montrer que la suite $S_{n}=\sum_{k=1}^{n}f(k) $ a une limite quand
$n\rightarrow \infty $ si et seulement si $\int_{1}^{\infty }f $ converge, et
que dans ce cas :
$$\int_{n+1}^{\infty }f\leq \lim_{m\rightarrow \infty
}\sum_{k=n+1}^{m}f(k)\leq \int_{n}^{\infty }f. $$}
    \item \question{Montrer que si $\int_{1}^{\infty }f $ diverge on a: $S_{n}\backsim
\int_{1}^{n}f $ quand $n\rightarrow \infty .$}
\end{enumerate}
}
